\documentclass[UTF8]{ctexart}
\usepackage[a4paper,margin=1.8cm]{geometry}
\usepackage{graphicx}
\usepackage{float}
\usepackage{longtable}
\usepackage{array}
\usepackage{xcolor}
\usepackage{fancyvrb}
\usepackage{hyperref}
\setlength{\parindent}{0pt}
\setlength{\parskip}{0.45em}
\begin{document}
\section*{补充模拟卷 03：中间代码表示补缺}

\subsection*{A 卷题目}

\subsubsection*{一、四元式、三元式、间接三元式（25 分）}

将表达式：

\begin{Verbatim}[fontsize=\small,breaklines=true]
x = (a + b) * (c - d) + e
\end{Verbatim}

翻译为：

1. 普通三地址码。 2. 四元式。 3. 三元式。 4. 间接三元式。

\subsubsection*{二、函数调用与参数传递（15 分）}

将下面语句翻译为三地址码：

\begin{Verbatim}[fontsize=\small,breaklines=true]
y = f(a, b + c)
\end{Verbatim}

要求体现 \texttt{param} 和 \texttt{call}。

\subsubsection*{三、数组访问与地址计算（20 分）}

假设 \texttt{int} 宽度为 4，将下面语句翻译为三地址码：

\begin{Verbatim}[fontsize=\small,breaklines=true]
val = a[i]
a[i] = y
\end{Verbatim}

要求体现下标地址计算。

\subsubsection*{四、fall-through 与减少冗余跳转（20 分）}

为下面语句生成三地址码，尽量利用顺序执行减少冗余 \texttt{goto}：

\begin{Verbatim}[fontsize=\small,breaklines=true]
if (x < 100 || x > 200 && x != y)
  x = 0;
\end{Verbatim}

说明 \texttt{fall-through} 的含义。

\subsubsection*{五、do-while、break、continue（20 分）}

1. 为下面语句生成三地址码：

\begin{Verbatim}[fontsize=\small,breaklines=true]
do {
  a = a + 1;
} while (a < 10 && b > 0);
\end{Verbatim}

2. 说明 \texttt{break} 和 \texttt{continue} 在循环中的跳转目标。

\subsection*{B 按考试标准重写答案}

\subsubsection*{一、四元式、三元式、间接三元式答案}

\paragraph*{1. 普通三地址码}

表达式：

\begin{Verbatim}[fontsize=\small,breaklines=true]
x = (a + b) * (c - d) + e
\end{Verbatim}

先计算括号内的两个子表达式，再乘法，再加 \texttt{e}，最后赋给 \texttt{x}：

\begin{longtable}{|p{0.31\linewidth}|p{0.31\linewidth}|p{0.31\linewidth}|}
\hline
步骤 & 子表达式 & 临时变量 \\
\hline
\endfirsthead
步骤 & 子表达式 & 临时变量 \\
\hline
\endhead
1 & \texttt{a + b} & \texttt{t1} \\
\hline
2 & \texttt{c - d} & \texttt{t2} \\
\hline
3 & \texttt{t1 * t2} & \texttt{t3} \\
\hline
4 & \texttt{t3 + e} & \texttt{t4} \\
\hline
5 & 赋值给 \texttt{x} & \texttt{x} \\
\hline
\end{longtable}

三地址码：

\begin{Verbatim}[fontsize=\small,breaklines=true]
t1 = a + b
t2 = c - d
t3 = t1 * t2
t4 = t3 + e
x = t4
\end{Verbatim}

\paragraph*{2. 四元式}

四元式格式为：

\begin{Verbatim}[fontsize=\small,breaklines=true]
(op, arg1, arg2, result)
\end{Verbatim}

\begin{longtable}{|p{0.19\linewidth}|p{0.19\linewidth}|p{0.19\linewidth}|p{0.19\linewidth}|p{0.19\linewidth}|}
\hline
序号 & op & arg1 & arg2 & result \\
\hline
\endfirsthead
序号 & op & arg1 & arg2 & result \\
\hline
\endhead
1 & \texttt{+} & \texttt{a} & \texttt{b} & \texttt{t1} \\
\hline
2 & \texttt{-} & \texttt{c} & \texttt{d} & \texttt{t2} \\
\hline
3 & \texttt{*} & \texttt{t1} & \texttt{t2} & \texttt{t3} \\
\hline
4 & \texttt{+} & \texttt{t3} & \texttt{e} & \texttt{t4} \\
\hline
5 & \texttt{=} & \texttt{t4} & \texttt{\textbackslash{}\_} & \texttt{x} \\
\hline
\end{longtable}

\paragraph*{3. 三元式}

三元式格式为：

\begin{Verbatim}[fontsize=\small,breaklines=true]
(op, arg1, arg2)
\end{Verbatim}

结果不再显式命名为 \texttt{t1}、\texttt{t2}，而是用三元式序号引用。

\begin{longtable}{|p{0.19\linewidth}|p{0.19\linewidth}|p{0.19\linewidth}|p{0.19\linewidth}|p{0.19\linewidth}|}
\hline
序号 & op & arg1 & arg2 & 含义 \\
\hline
\endfirsthead
序号 & op & arg1 & arg2 & 含义 \\
\hline
\endhead
0 & \texttt{+} & \texttt{a} & \texttt{b} & \texttt{(a+b)} \\
\hline
1 & \texttt{-} & \texttt{c} & \texttt{d} & \texttt{(c-d)} \\
\hline
2 & \texttt{*} & \texttt{(0)} & \texttt{(1)} & \texttt{(a+b)*(c-d)} \\
\hline
3 & \texttt{+} & \texttt{(2)} & \texttt{e} & \texttt{(a+b)*(c-d)+e} \\
\hline
4 & \texttt{=} & \texttt{(3)} & \texttt{x} & \texttt{x=(3)} \\
\hline
\end{longtable}

写成列表：

\begin{Verbatim}[fontsize=\small,breaklines=true]
0: (+, a, b)
1: (-, c, d)
2: (*, (0), (1))
3: (+, (2), e)
4: (=, (3), x)
\end{Verbatim}

\paragraph*{4. 间接三元式}

间接三元式由“三元式表”和“指针表”组成。三元式表保持不变：

\begin{longtable}{|p{0.23\linewidth}|p{0.23\linewidth}|p{0.23\linewidth}|p{0.23\linewidth}|}
\hline
三元式序号 & op & arg1 & arg2 \\
\hline
\endfirsthead
三元式序号 & op & arg1 & arg2 \\
\hline
\endhead
0 & \texttt{+} & \texttt{a} & \texttt{b} \\
\hline
1 & \texttt{-} & \texttt{c} & \texttt{d} \\
\hline
2 & \texttt{*} & \texttt{(0)} & \texttt{(1)} \\
\hline
3 & \texttt{+} & \texttt{(2)} & \texttt{e} \\
\hline
4 & \texttt{=} & \texttt{(3)} & \texttt{x} \\
\hline
\end{longtable}

指针表决定实际执行顺序：

\begin{longtable}{|p{0.31\linewidth}|p{0.31\linewidth}|p{0.31\linewidth}|}
\hline
指针项 & 指向三元式 & 当前执行内容 \\
\hline
\endfirsthead
指针项 & 指向三元式 & 当前执行内容 \\
\hline
\endhead
p0 & 0 & \texttt{a+b} \\
\hline
p1 & 1 & \texttt{c-d} \\
\hline
p2 & 2 & \texttt{(0)*(1)} \\
\hline
p3 & 3 & \texttt{(2)+e} \\
\hline
p4 & 4 & \texttt{x=(3)} \\
\hline
\end{longtable}

间接三元式的优点是：优化时若要调整执行顺序，可主要调整指针表，而不必移动三元式本身。

\subsubsection*{二、函数调用与参数传递答案}

语句：

\begin{Verbatim}[fontsize=\small,breaklines=true]
y = f(a, b + c)
\end{Verbatim}

先计算非简单实参 \texttt{b + c}，再传参并调用函数。三地址码为：

\begin{Verbatim}[fontsize=\small,breaklines=true]
t1 = b + c
param a
param t1
t2 = call f, 2
y = t2
\end{Verbatim}

逐行说明：

\begin{longtable}{|p{0.47\linewidth}|p{0.47\linewidth}|}
\hline
行 & 作用 \\
\hline
\endfirsthead
行 & 作用 \\
\hline
\endhead
\texttt{t1 = b + c} & 先求第二个实参表达式 \\
\hline
\texttt{param a} & 传入第一个实参 \\
\hline
\texttt{param t1} & 传入第二个实参 \\
\hline
\texttt{t2 = call f, 2} & 调用 \texttt{f}，实参数量为 2，返回值放入 \texttt{t2} \\
\hline
\texttt{y = t2} & 把返回值赋给 \texttt{y} \\
\hline
\end{longtable}

若课堂约定参数从右向左压栈，则 \texttt{param} 顺序可按约定调整；本答案采用题面常见的从左到右传参写法。

\subsubsection*{三、数组访问与地址计算答案}

题目给定 \texttt{int} 宽度为 4。若数组 \texttt{a} 的基地址记为 \texttt{a}，则：

\begin{Verbatim}[fontsize=\small,breaklines=true]
a[i] 的地址 = a + i * 4
\end{Verbatim}

\paragraph*{1. `val = a[i]`}

\begin{longtable}{|p{0.31\linewidth}|p{0.31\linewidth}|p{0.31\linewidth}|}
\hline
步骤 & 计算内容 & 三地址码 \\
\hline
\endfirsthead
步骤 & 计算内容 & 三地址码 \\
\hline
\endhead
1 & 计算字节偏移 & \texttt{t1 = i * 4} \\
\hline
2 & 计算元素地址 & \texttt{t2 = a + t1} \\
\hline
3 & 从地址取值 & \texttt{val = *t2} \\
\hline
\end{longtable}

完整代码：

\begin{Verbatim}[fontsize=\small,breaklines=true]
t1 = i * 4
t2 = a + t1
val = *t2
\end{Verbatim}

\paragraph*{2. `a[i] = y`}

\begin{longtable}{|p{0.31\linewidth}|p{0.31\linewidth}|p{0.31\linewidth}|}
\hline
步骤 & 计算内容 & 三地址码 \\
\hline
\endfirsthead
步骤 & 计算内容 & 三地址码 \\
\hline
\endhead
1 & 计算字节偏移 & \texttt{t3 = i * 4} \\
\hline
2 & 计算元素地址 & \texttt{t4 = a + t3} \\
\hline
3 & 写入地址 & \texttt{*t4 = y} \\
\hline
\end{longtable}

完整代码：

\begin{Verbatim}[fontsize=\small,breaklines=true]
t3 = i * 4
t4 = a + t3
*t4 = y
\end{Verbatim}

注意：\texttt{val = a[i]} 是取右值，需要从地址读；\texttt{a[i] = y} 是写左值，需要先算出地址再写入。

\subsubsection*{四、fall-through 与减少冗余跳转答案}

\paragraph*{1. 条件结构分析}

语句：

\begin{Verbatim}[fontsize=\small,breaklines=true]
if (x < 100 || x > 200 && x != y)
  x = 0;
\end{Verbatim}

\texttt{\textbackslash{}\&\textbackslash{}\&} 优先级高于 \texttt{||}，所以条件等价于：

\begin{Verbatim}[fontsize=\small,breaklines=true]
(x < 100) || ((x > 200) && (x != y))
\end{Verbatim}

短路逻辑如下：

\begin{longtable}{|p{0.31\linewidth}|p{0.31\linewidth}|p{0.31\linewidth}|}
\hline
子条件 & 为真时 & 为假时 \\
\hline
\endfirsthead
子条件 & 为真时 & 为假时 \\
\hline
\endhead
\texttt{x < 100} & 整体为真，执行 then & 继续判断 \texttt{x > 200} \\
\hline
\texttt{x > 200} & 继续判断 \texttt{x != y} & 整体为假，跳到出口 \\
\hline
\texttt{x != y} & 整体为真，执行 then & 整体为假，跳到出口 \\
\hline
\end{longtable}

\paragraph*{2. 利用 fall-through 的三地址码}

把 then 分支放在条件判断之后，让最后一个真路径自然落入 then 分支：

\begin{Verbatim}[fontsize=\small,breaklines=true]
if x < 100 goto L_then
if x <= 200 goto L_next
if x == y goto L_next

L_then:
x = 0

L_next:
\end{Verbatim}

逐行解释：

\begin{longtable}{|p{0.23\linewidth}|p{0.23\linewidth}|p{0.23\linewidth}|p{0.23\linewidth}|}
\hline
代码 & 说明 &  &  \\
\hline
\endfirsthead
代码 & 说明 &  &  \\
\hline
\endhead
\texttt{if x < 100 goto L\textbackslash{}\_then} & 左侧 ` &  & ` 为真，直接执行 then \\
\hline
\texttt{if x <= 200 goto L\textbackslash{}\_next} & \texttt{x > 200} 为假，整体条件为假 &  &  \\
\hline
\texttt{if x == y goto L\textbackslash{}\_next} & \texttt{x != y} 为假，整体条件为假 &  &  \\
\hline
顺序落入 \texttt{L\textbackslash{}\_then} & 说明 \texttt{x > 200} 且 \texttt{x != y} 都为真 &  &  \\
\hline
\end{longtable}

\paragraph*{3. fall-through 含义}

\texttt{fall-through} 指“不生成显式跳转指令，而让控制流按代码顺序自然进入下一条指令”。本题中第三个判断之后没有写：

\begin{Verbatim}[fontsize=\small,breaklines=true]
goto L_then
\end{Verbatim}

而是让满足条件的路径直接顺序执行到 \texttt{L\textbackslash{}\_then}，减少了一条冗余跳转。

\subsubsection*{五、do-while、break、continue 答案}

\paragraph*{1. do-while 三地址码}

源程序：

\begin{Verbatim}[fontsize=\small,breaklines=true]
do {
  a = a + 1;
} while (a < 10 && b > 0);
\end{Verbatim}

\texttt{do-while} 是后测试循环，循环体至少执行一次。条件是：

\begin{Verbatim}[fontsize=\small,breaklines=true]
(a < 10) && (b > 0)
\end{Verbatim}

\texttt{\textbackslash{}\&\textbackslash{}\&} 的短路规则是：第一个条件为假时，不再检查第二个条件，直接退出循环。

三地址码：

\begin{Verbatim}[fontsize=\small,breaklines=true]
L_body:
t1 = a + 1
a = t1

L_test:
if a < 10 goto L_check_b
goto L_exit

L_check_b:
if b > 0 goto L_body
goto L_exit

L_exit:
\end{Verbatim}

流程说明：

\begin{longtable}{|p{0.47\linewidth}|p{0.47\linewidth}|}
\hline
标签 & 作用 \\
\hline
\endfirsthead
标签 & 作用 \\
\hline
\endhead
\texttt{L\textbackslash{}\_body} & 循环体开始位置 \\
\hline
\texttt{L\textbackslash{}\_test} & do-while 底部条件检查位置 \\
\hline
\texttt{L\textbackslash{}\_check\textbackslash{}\_b} & 只有 \texttt{a < 10} 为真时才检查 \texttt{b > 0} \\
\hline
\texttt{L\textbackslash{}\_exit} & 循环出口 \\
\hline
\end{longtable}

\paragraph*{2. `break` 和 `continue` 的跳转目标}

在循环翻译中，\texttt{break} 和 \texttt{continue} 不是普通顺序执行语句，它们直接改变控制流：

\begin{longtable}{|p{0.31\linewidth}|p{0.31\linewidth}|p{0.31\linewidth}|}
\hline
语句 & while 循环中的目标 & do-while 循环中的目标 \\
\hline
\endfirsthead
语句 & while 循环中的目标 & do-while 循环中的目标 \\
\hline
\endhead
\texttt{break} & 跳到循环出口 & 跳到循环出口 \\
\hline
\texttt{continue} & 跳到循环头的条件判断 & 跳到底部条件判断 \\
\hline
\end{longtable}

对应本题的 \texttt{do-while}：

\begin{Verbatim}[fontsize=\small,breaklines=true]
break    -> goto L_exit
continue -> goto L_test
\end{Verbatim}

原因是 \texttt{continue} 表示结束本轮循环并进入下一轮判断；对 \texttt{do-while}，下一轮判断发生在循环体之后的测试位置。
\end{document}
